\documentclass[11pt]{article}
\usepackage[margin=1in]{geometry}
\usepackage{amsmath,amssymb,amsthm,mathtools}
\usepackage[T1]{fontenc}
\usepackage[utf8]{inputenc}
\usepackage{enumitem}
\usepackage{booktabs,longtable,array}
\usepackage{hyperref}
\usepackage{xcolor}
\usepackage{microtype}
\hypersetup{colorlinks=true, linkcolor=blue!50!black, citecolor=blue!50!black, urlcolor=blue!50!black}

\newtheorem{theorem}{Theorem}[section]
\newtheorem{proposition}[theorem]{Proposition}
\newtheorem{lemma}[theorem]{Lemma}
\newtheorem{corollary}[theorem]{Corollary}
\newtheorem{claim}[theorem]{Claim}
\theoremstyle{definition}
\newtheorem{definition}[theorem]{Definition}
\newtheorem{assumption}[theorem]{Assumption}
\theoremstyle{remark}
\newtheorem{remark}[theorem]{Remark}

\newcommand{\KL}{\mathrm{KL}}
\newcommand{\Ent}{\mathrm{Ent}}
\newcommand{\Prob}{\mathsf{P}}
\newcommand{\E}{\mathbb{E}}
\newcommand{\R}{\mathbb{R}}
\newcommand{\N}{\mathbb{N}}
\newcommand{\Pcal}{\mathcal{P}}
\newcommand{\Mcal}{\mathcal{M}}
\newcommand{\Qcal}{\mathcal{Q}}
\newcommand{\Zcal}{\mathcal{Z}}
\newcommand{\Fcal}{\mathcal{F}}
\newcommand{\Gcal}{\mathcal{G}}
\newcommand{\W}{\mathcal{W}}
\newcommand{\one}{\mathbf{1}}
\newcommand{\dd}{\,d}
\newcommand{\supp}{\operatorname{supp}}
\newcommand{\argmin}{\operatorname*{arg\,min}}
\newcommand{\argmax}{\operatorname*{arg\,max}}
\newcommand{\esssup}{\operatorname*{ess\,sup}}
\newcommand{\ri}{\operatorname{ri}}
\newcommand{\dom}{\operatorname{dom}}
\newcommand{\Law}{\operatorname{Law}}
\newcommand{\Var}{\operatorname{Var}}
\newcommand{\Cov}{\operatorname{Cov}}
\newcommand{\diag}{\operatorname{diag}}
\newcommand{\Id}{\mathrm{Id}}
\newcommand{\Lip}{\operatorname{Lip}}
\newcommand{\osc}{\operatorname{osc}}
\newcommand{\diam}{\operatorname{diam}}
\newcommand{\proj}{\operatorname{proj}}
\newcommand{\e}{\mathrm{e}}

\newenvironment{argument}{\paragraph{Argument.}\begin{enumerate}[leftmargin=2em]}{\end{enumerate}}
\newenvironment{proofoutline}{\paragraph{Proof architecture.}\begin{enumerate}[leftmargin=2em]}{\end{enumerate}}


% Source-faithfulness apparatus used by the paper reconstructions.
\newcommand{\sourceanchor}[1]{\par\smallskip\noindent\textbf{Source anchor.} #1\par\smallskip}
\newcommand{\sourcecomment}[1]{\begin{remark}#1\end{remark}}
\newenvironment{sourceguide}
  {\begin{description}[style=nextline,leftmargin=3.2cm,labelwidth=3cm]}
  {\end{description}}

% Backward-compatible environments used by some of the older reconstructions.
\newenvironment{distilled}[1][]{\begin{theorem}[#1]}{\end{theorem}}
\newenvironment{distillation}{\begin{enumerate}[leftmargin=2em]}{\end{enumerate}}

\newenvironment{commentarynotes}{\begin{remark}\begin{enumerate}[leftmargin=2em]}{\end{enumerate}\end{remark}}

\title{Data-Driven Distributionally Robust Optimization with the Wasserstein Metric}
\author{}
\date{Mohajerin Esfahani and Kuhn (2018)}
\begin{document}
\maketitle

\section*{Source map and scope}
\begin{sourceguide}
\item[Primary source] Peyman Mohajerin Esfahani and Daniel Kuhn, \emph{Data-Driven Distributionally Robust Optimization Using the Wasserstein Metric: Performance Guarantees and Tractable Reformulations} \cite{MohajerinEsfahaniKuhn2018}.
\item[Coverage] The mathematical chain in Sections 3--4 and Appendix A that supports the empirical Wasserstein ambiguity set: Definition 3.1, Theorem 3.2, Assumption 3.3, Theorems 3.4--3.6, Lemma 3.7, Assumption 4.1, Theorem 4.2, Corollary 4.3, Theorem 4.4, Lemma 4.5, and the upper-Lipschitz-majorant lemma used in asymptotic consistency.
\item[Target results] The finite-sample guarantee, asymptotic consistency, exact finite convex reformulation of the empirical worst-case expectation, and finite-atomic extremal-law construction.
\item[Method] The proof retains the source's generalized moment problem, Lagrangian dual, robust-counterpart, minimax, conjugate-infimal-convolution, and perspective-duality route.  It does not substitute a later general optimal-transport duality theorem.
\end{sourceguide}

\section{Empirical Wasserstein ambiguity sets}
Let $\Xi\subseteq\R^m$ be closed, let $\|\cdot\|$ be a norm on $\R^m$, and let $\Mcal(\Xi)$ denote the probability measures on $\Xi$ having finite first moment.  For $Q_1,Q_2\in\Mcal(\Xi)$, define
\begin{equation}\label{eq:MEK-W1}
 d_W(Q_1,Q_2)
 :=\inf_{\Pi\in\Pi(Q_1,Q_2)}
   \int_{\Xi\times\Xi}\|\xi_1-\xi_2\|\,\Pi(d\xi_1,d\xi_2).
\end{equation}
For independent samples $\widehat\xi_1,\dots,\widehat\xi_N$ from an unknown law $P$, write
\[
 \widehat P_N:=\frac1N\sum_{i=1}^N\delta_{\widehat\xi_i},
 \qquad
 B_\varepsilon(\widehat P_N)
 :=\{Q\in\Mcal(\Xi):d_W(Q,\widehat P_N)\le\varepsilon\}.
\]

The Kantorovich--Rubinstein representation used in the statistical argument is
\begin{equation}\label{eq:MEK-KR}
 d_W(Q_1,Q_2)
 =\sup_{f:\,\Lip(f)\le1}
 \left\{\int f\,dQ_1-\int f\,dQ_2\right\}.
\end{equation}
The paper cites the classical theorem and a modern unbounded-support treatment rather than proving \eqref{eq:MEK-KR}; the reconstruction uses it exactly where the source does.

For a decision $x\in\mathcal X$ and loss $h(x,\xi)$, the distributionally robust problem is
\begin{equation}\label{eq:MEK-DRO}
 \widehat J_N
 :=\inf_{x\in\mathcal X}
 \sup_{Q\in B_{\varepsilon_N}(\widehat P_N)}
 \E_Q[h(x,\xi)].
\end{equation}
Let $J^*:=\inf_{x\in\mathcal X}\E_P[h(x,\xi)]$ denote the true stochastic-program value.

\section{Measure concentration and the radius choice}
\sourceanchor{Assumption 3.3 and Theorem 3.4.}
Assume that there is $a>1$ such that
\begin{equation}\label{eq:MEK-light-tail}
 A:=\E_P[\exp(\|\xi\|^a)]<\infty.
\end{equation}
The paper imports the following empirical-measure concentration estimate from Fournier and Guillin.

\begin{theorem}[Source Theorem 3.4: concentration]\label{thm:MEK-concentration}
There are $c_1,c_2>0$, depending only on $a$, $A$, and $m$, such that for $N\ge1$, $m\ne2$, and $\varepsilon>0$,
\begin{equation}\label{eq:MEK-concentration}
 P^N\{d_W(P,\widehat P_N)\ge\varepsilon\}
 \le
 \begin{cases}
 c_1\exp[-c_2N\varepsilon^{\max\{m,2\}}],&\varepsilon\le1,\\
 c_1\exp[-c_2N\varepsilon^a],&\varepsilon>1.
 \end{cases}
\end{equation}
\end{theorem}

For confidence level $1-\beta$, invert the right-hand side:
\begin{equation}\label{eq:MEK-radius}
 \varepsilon_N(\beta):=
 \begin{cases}
 \left(\dfrac{\log(c_1/\beta)}{c_2N}\right)^{1/\max\{m,2\}},
 &N\ge\dfrac{\log(c_1/\beta)}{c_2},\\[1.2ex]
 \left(\dfrac{\log(c_1/\beta)}{c_2N}\right)^{1/a},
 &N<\dfrac{\log(c_1/\beta)}{c_2}.
 \end{cases}
\end{equation}
Then
\[
 P^N\{P\in B_{\varepsilon_N(\beta)}(\widehat P_N)\}\ge1-\beta.
\]

\begin{theorem}[Source Theorem 3.5: finite-sample guarantee]\label{thm:MEK-finite-sample}
Let $\widehat x_N$ solve \eqref{eq:MEK-DRO} with radius \eqref{eq:MEK-radius}.  Then
\begin{equation}\label{eq:MEK-finite-guarantee}
 P^N\left\{
 \E_P[h(\widehat x_N,\xi)]\le\widehat J_N
 \right\}\ge1-\beta.
\end{equation}
\end{theorem}
\begin{proof}
On the event $P\in B_{\varepsilon_N(\beta)}(\widehat P_N)$,
\[
 \E_P[h(\widehat x_N,\xi)]
 \le\sup_{Q\in B_{\varepsilon_N(\beta)}(\widehat P_N)}
 \E_Q[h(\widehat x_N,\xi)]
 =\widehat J_N.
\]
The concentration theorem bounds the complement of this event by $\beta$.
\end{proof}

\section{Asymptotic consistency}
\sourceanchor{Theorem 3.6, Lemma 3.7, and Appendix Lemma A.1.}
Let $\beta_N\in(0,1)$ satisfy
\begin{equation}\label{eq:MEK-beta-summable}
 \sum_{N=1}^\infty\beta_N<\infty,
 \qquad
 \varepsilon_N(\beta_N)\longrightarrow0.
\end{equation}

\begin{lemma}[Source Lemma 3.7]\label{lem:MEK-distribution-convergence}
For any possibly data-dependent sequence
\[
 \widehat Q_N\in B_{\varepsilon_N(\beta_N)}(\widehat P_N),
\]
one has
\[
 d_W(P,\widehat Q_N)\longrightarrow0
 \quad P^\infty\text{-almost surely}.
\]
\end{lemma}
\begin{proof}
The triangle inequality gives
\[
 d_W(P,\widehat Q_N)
 \le d_W(P,\widehat P_N)+\varepsilon_N(\beta_N).
\]
Theorem~\ref{thm:MEK-concentration} and the radius definition imply
\[
 P^N\{d_W(P,\widehat P_N)>\varepsilon_N(\beta_N)\}\le\beta_N.
\]
By Borel--Cantelli, almost surely the first term is at most $\varepsilon_N(\beta_N)$ for all sufficiently large $N$, and the sum tends to zero.
\end{proof}

The proof of value convergence uses a decreasing Lipschitz approximation from above.

\begin{lemma}[Source Appendix Lemma A.1]\label{lem:MEK-lipschitz-majorants}
Let $g:\Xi\to\R$ be upper semicontinuous and satisfy
\[
 |g(\xi)|\le L(1+\|\xi\|).
\]
For $k\ge2$, define
\begin{equation}\label{eq:MEK-majorant}
 g_k(\xi):=\sup_{\xi'\in\Xi}
 \{g(\xi')-kL\|\xi-\xi'\|\}.
\end{equation}
Then $g_k$ is $kL$-Lipschitz, $g_k\ge g$, $g_k\downarrow g$ pointwise, and $g_k(\xi)\le L(1+\|\xi\|)$ up to an additive constant independent of $\xi$.
\end{lemma}
\begin{proof}
Each function $\xi\mapsto g(\xi')-kL\|\xi-\xi'\|$ is $kL$-Lipschitz, so the supremum is too.  Choosing $\xi'=\xi$ gives $g_k\ge g$, and increasing $k$ decreases the supremum.  If $\xi_k'$ is a $1/k$-maximizer, linear growth gives
\[
 g(\xi)
 \le g_k(\xi)
 \le L(1+\|\xi_k'\|)-kL\|\xi-\xi_k'\|+\frac1k.
\]
Using $\|\xi_k'\|\le\|\xi\|+\|\xi_k'-\xi\|$ shows $\xi_k'\to\xi$.  Upper semicontinuity then gives $\limsup g_k(\xi)\le g(\xi)$.
\end{proof}

\begin{theorem}[Source Theorem 3.6]\label{thm:MEK-consistency}
Assume \eqref{eq:MEK-light-tail} and \eqref{eq:MEK-beta-summable}.
\begin{enumerate}[label=(\roman*),leftmargin=2em]
\item If $h(x,\cdot)$ is upper semicontinuous and
\[
 |h(x,\xi)|\le L(1+\|\xi\|)
\]
uniformly in $x$, then $\widehat J_N\to J^*$ almost surely.
\item If additionally $\mathcal X$ is closed and $h(\cdot,\xi)$ is lower semicontinuous for every $\xi$, then every accumulation point of $\widehat x_N$ is almost surely optimal for the true problem.
\end{enumerate}
\end{theorem}
\begin{proof}
The finite-sample theorem and Borel--Cantelli imply that, almost surely, for all large $N$,
\begin{equation}\label{eq:MEK-eventual-sandwich}
 J^*\le\E_P[h(\widehat x_N,\xi)]\le\widehat J_N.
\end{equation}
Fix $\delta>0$ and choose $x_\delta$ with
\[
 \E_P[h(x_\delta,\xi)]\le J^*+\delta.
\]
Choose $\widehat Q_N$ in the ambiguity ball that is $\delta$-optimal for the inner supremum at $x_\delta$.  Apply Lemma~\ref{lem:MEK-lipschitz-majorants} to $h(x_\delta,\cdot)$ and use \eqref{eq:MEK-KR}:
\begin{align*}
 \limsup_N\widehat J_N
 &\le\limsup_N\E_{\widehat Q_N}[h(x_\delta,\xi)]+\delta\\
 &\le\lim_{k\to\infty}\limsup_N
 \left(\E_P[h_k(x_\delta,\xi)]
 +L_kd_W(P,\widehat Q_N)\right)+\delta\\
 &=\E_P[h(x_\delta,\xi)]+\delta
 \le J^*+2\delta.
\end{align*}
Here Lemma~\ref{lem:MEK-distribution-convergence} removes the Wasserstein term, and monotone convergence from above is justified by the common integrable linear-growth bound.  Letting $\delta\downarrow0$ and combining with \eqref{eq:MEK-eventual-sandwich} proves (i).

For (ii), take a convergent subsequence $\widehat x_{N_j}\to x^*$.  Closedness gives $x^*\in\mathcal X$.  Lower semicontinuity and Fatou's lemma yield
\[
 J^*\le\E_P[h(x^*,\xi)]
 \le\liminf_j\E_P[h(\widehat x_{N_j},\xi)]
 \le\lim_j\widehat J_{N_j}=J^*.
\]
Thus $x^*$ is optimal.
\end{proof}

The paper supplies counterexamples showing that upper semicontinuity in $\xi$, linear growth, and lower semicontinuity in $x$ cannot simply be dropped: a jump loss at the support point, a quadratic loss under a first-Wasserstein ball, and a discontinuous decision loss respectively break the conclusions.

\section{The empirical worst-case expectation}
Suppress the decision variable and write
\begin{equation}\label{eq:MEK-wce}
 \mathcal V_\varepsilon(\ell)
 :=\sup_{Q\in B_\varepsilon(\widehat P_N)}\E_Q[\ell(\xi)].
\end{equation}
Assume
\begin{equation}\label{eq:MEK-piecewise}
 \ell(\xi)=\max_{1\le k\le K}\ell_k(\xi).
\end{equation}

\begin{assumption}[Source Assumption 4.1]\label{ass:MEK-convexity}
The set $\Xi$ is closed and convex.  Each $-\ell_k$ is proper, convex, and lower semicontinuous, and no $\ell_k$ is identically $-\infty$ on $\Xi$.
\end{assumption}

Let $\chi_\Xi$ denote the convex indicator of $\Xi$, $\sigma_\Xi=\chi_\Xi^*$ its support function, and $\|\cdot\|_*$ the dual norm.

\section{Disintegration into a generalized moment problem}
Every coupling of $Q$ with the empirical law has the form
\[
 \Pi=\frac1N\sum_{i=1}^N\delta_{\widehat\xi_i}\otimes Q_i
\]
for conditional laws $Q_i\in\Mcal(\Xi)$.  Hence
\begin{equation}\label{eq:MEK-moment-primal}
 \mathcal V_\varepsilon(\ell)
 =\sup_{Q_1,\dots,Q_N}
 \left\{
 \frac1N\sum_{i=1}^N\int_\Xi\ell(\xi)Q_i(d\xi):
 \frac1N\sum_{i=1}^N\int_\Xi\|\xi-\widehat\xi_i\|Q_i(d\xi)
 \le\varepsilon
 \right\}.
\end{equation}
Dualize the single moment constraint with $\lambda\ge0$.  Weak max--min duality gives
\begin{align}
 \mathcal V_\varepsilon(\ell)
 &\le\inf_{\lambda\ge0}
 \left\{
 \lambda\varepsilon
 +\frac1N\sum_{i=1}^N
 \sup_{\xi\in\Xi}
 [\ell(\xi)-\lambda\|\xi-\widehat\xi_i\|]
 \right\}.
 \label{eq:MEK-scalar-dual}
\end{align}
Under Assumption~\ref{ass:MEK-convexity}, the generalized moment-problem duality cited by the paper makes this an equality for $\varepsilon>0$.  At $\varepsilon=0$, the multiplier can diverge without objective penalty; the source checks directly that the right-hand side decreases to the sample average.

\section{Robust-counterpart derivation}
Introduce epigraph variables $s_i$ in \eqref{eq:MEK-scalar-dual}:
\begin{equation}\label{eq:MEK-semi-infinite}
 \sup_{\xi\in\Xi}
 [\ell_k(\xi)-\lambda\|\xi-\widehat\xi_i\|]
 \le s_i,
 \qquad i\le N,
 \quad k\le K.
\end{equation}
The dual-norm formula gives
\[
 \lambda\|\xi-\widehat\xi_i\|
 =\max_{\|z\|_*\le\lambda}
 \langle z,\xi-\widehat\xi_i\rangle.
\]
Therefore
\begin{align*}
 &\sup_{\xi\in\Xi}
 \left(\ell_k(\xi)-
 \max_{\|z\|_*\le\lambda}
 \langle z,\xi-\widehat\xi_i\rangle\right)\\
 &\qquad\le
 \min_{\|z_{ik}\|_*\le\lambda}
 \sup_{\xi\in\Xi}
 [\ell_k(\xi)-\langle z_{ik},\xi\rangle
 +\langle z_{ik},\widehat\xi_i\rangle].
\end{align*}
Without convexity this is only an upper bound.  Under Assumption~\ref{ass:MEK-convexity}, the compactness of the dual-norm ball and the source's minimax theorem make it an equality.

The remaining supremum is a conjugate:
\[
 \sup_{\xi\in\Xi}[\ell_k(\xi)-\langle z,\xi\rangle]
 =[-\ell_k+\chi_\Xi]^*(-z).
\]
After replacing $-z$ by $z$, the sum-conjugate formula gives
\begin{equation}\label{eq:MEK-inf-conv}
 [-\ell_k+\chi_\Xi]^*(z)
 =\operatorname{cl}\inf_{\nu\in\R^m}
 \bigl\{[-\ell_k]^*(z-\nu)+\sigma_\Xi(\nu)\bigr\}.
\end{equation}
The closure does not alter a sublevel condition of the form ``$\le s_i$,'' so the auxiliary variable $\nu_{ik}$ yields the finite program.

\begin{theorem}[Source Theorem 4.2]\label{thm:MEK-convex-reduction}
Under Assumption~\ref{ass:MEK-convexity}, \eqref{eq:MEK-wce} equals
\begin{equation}\label{eq:MEK-finite-program}
\begin{aligned}
 \inf_{\lambda,s,z,\nu}\quad
 &\lambda\varepsilon+\frac1N\sum_{i=1}^Ns_i\\
 \text{subject to}\quad
 &[-\ell_k]^*(z_{ik}-\nu_{ik})
 +\sigma_\Xi(\nu_{ik})
 -\langle z_{ik},\widehat\xi_i\rangle
 \le s_i,
 &&i\le N,\ k\le K,\\
 &\|z_{ik}\|_*\le\lambda,
 &&i\le N,\ k\le K,\\
 &\lambda\ge0.
\end{aligned}
\end{equation}
\end{theorem}

\begin{corollary}[Source Corollary 4.3]\label{cor:MEK-upper-reduction}
Without Assumption~\ref{ass:MEK-convexity}, the pre-infimal-convolution finite program remains a conservative upper bound on \eqref{eq:MEK-wce}.
\end{corollary}
The distinction is important: the max--min and minimax exchanges are exact only under the source's duality and convexity hypotheses.

\section{Extremal distributions by dualizing the finite program}
\sourceanchor{Theorem 4.4 and Lemma 4.5.}
The source dualizes the finite convex program rather than invoking an abstract finite-support theorem.

\begin{lemma}[Source Lemma 4.5: perspective identity]\label{lem:MEK-perspective}
Let $f:\R^m\to\overline\R$ be proper, convex, and lower semicontinuous, and let $\widehat\xi\in\R^m$.  Define
\[
 F(q,\alpha)
 :=\inf_{z\in\R^m}
 \{\langle z,q\rangle
 -\alpha\langle z,\widehat\xi\rangle
 +\alpha f^*(z)\},
 \qquad \alpha\ge0.
\]
Then
\begin{equation}\label{eq:MEK-perspective}
 F(q,\alpha)=
 \begin{cases}
 -\alpha f(\widehat\xi-q/\alpha),&\alpha>0,\\
 -\chi_{\{0\}}(q),&\alpha=0.
 \end{cases}
\end{equation}
\end{lemma}
\begin{proof}
For $\alpha=0$, the infimum of $\langle z,q\rangle$ over all $z$ is zero if $q=0$ and $-\infty$ otherwise.  For $\alpha>0$,
\[
 F(q,\alpha)
 =-[\alpha f^*]^*(\alpha\widehat\xi-q)
 =-\alpha f(\widehat\xi-q/\alpha)
\]
by Fenchel--Moreau and the scaling rule for conjugates.
\end{proof}

The Lagrange dual of the finite program introduces weights $\alpha_{ik}\ge0$ and transportation vectors $q_{ik}$.  Minimization over $s_i$ gives
\[
 \sum_{k=1}^K\alpha_{ik}=1
 \qquad(i=1,\dots,N),
\]
and minimization over $\lambda$ gives
\[
 \frac1N\sum_{i=1}^N\sum_{k=1}^K\|q_{ik}\|\le\varepsilon.
\]
Lemma~\ref{lem:MEK-perspective}, with $f=-\ell_k+\chi_\Xi$, gives the objective and support constraints.

\begin{theorem}[Source Theorem 4.4]\label{thm:MEK-worst-case-laws}
Under Assumption~\ref{ass:MEK-convexity}, \eqref{eq:MEK-wce} equals
\begin{equation}\label{eq:MEK-extremal-program}
\begin{aligned}
 \sup_{\alpha,q}\quad
 &\frac1N\sum_{i=1}^N\sum_{k=1}^K
 \alpha_{ik}\,
 \ell_k\!\left(\widehat\xi_i-\frac{q_{ik}}{\alpha_{ik}}\right)\\
 \text{subject to}\quad
 &\frac1N\sum_{i,k}\|q_{ik}\|\le\varepsilon,\\
 &\sum_{k=1}^K\alpha_{ik}=1,
 &&i=1,\dots,N,\\
 &\alpha_{ik}\ge0,
 &&i\le N,\ k\le K,\\
 &\widehat\xi_i-q_{ik}/\alpha_{ik}\in\Xi,
 &&i\le N,\ k\le K.
\end{aligned}
\end{equation}
The extended perspective convention is used: if $\alpha_{ik}=0$ and $q_{ik}\ne0$, the support constraint fails; if both vanish, the objective contribution is zero.

If $(\alpha(r),q(r))$ is a maximizing sequence for \eqref{eq:MEK-extremal-program}, set
\[
 \xi_{ik}(r):=\widehat\xi_i-\frac{q_{ik}(r)}{\alpha_{ik}(r)},
 \qquad
 Q_r:=\frac1N\sum_{i=1}^N\sum_{k=1}^K
 \alpha_{ik}(r)\delta_{\xi_{ik}(r)}.
\]
Then $Q_r\in B_\varepsilon(\widehat P_N)$ and
\[
 \E_{Q_r}[\ell]\longrightarrow\mathcal V_\varepsilon(\ell).
\]
If \eqref{eq:MEK-extremal-program} attains its optimum, the associated finite atomic measure is a worst-case distribution.
\end{theorem}

\begin{proof}
Strong duality for the finite program holds because the conjugate constraint has a free epigraph variable.  The perspective calculation gives \eqref{eq:MEK-extremal-program}.  For any feasible $(\alpha,q)$, transport mass $\alpha_{ik}/N$ from $\widehat\xi_i$ to $\xi_{ik}$.  The resulting coupling has cost
\[
 \frac1N\sum_{i,k}\alpha_{ik}
 \|\xi_{ik}-\widehat\xi_i\|
 =\frac1N\sum_{i,k}\|q_{ik}\|
 \le\varepsilon.
\]
Hence $Q\in B_\varepsilon(\widehat P_N)$.  Since $\ell=\max_k\ell_k$,
\[
 \E_Q[\ell]
 \ge\frac1N\sum_{i,k}
 \alpha_{ik}\ell_k(\xi_{ik}),
\]
which is the dual objective.  Equality of the finite primal and dual optimal values and weak duality for the original measure problem force convergence of the generated laws to the worst-case value.
\end{proof}

The construction uses at most $NK$ atoms.  It is not merely a generic Carath\'eodory reduction: each atom is tied to one empirical point and one active constituent of the pointwise-maximum loss.

\section*{Source issues and fidelity notes}
\begin{itemize}[leftmargin=2em]
\item The concentration theorem is imported from Fournier--Guillin; the paper does not reprove it.  The finite-sample theorem is the immediate confidence-set consequence reproduced above.
\item The asymptotic proof uses a decreasing sequence of Lipschitz majorants, not truncation of the loss or a generic epi-convergence theorem.
\item In Theorem 4.2, equality at $\varepsilon=0$ is handled separately because ordinary interior-point moment duality does not apply directly to a zero budget.
\item The passage from the semi-infinite constraint to the robust counterpart has two logically distinct steps: generalized-moment strong duality and the minimax interchange over the compact dual-norm ball.  Both are preserved.
\item The source's extended-arithmetic convention at $\alpha=0$ is part of Theorem 4.4.  Removing it would change the feasible set of the perspective program.
\item The finite-support laws in Theorem 4.4 are asymptotically worst-case for a maximizing sequence; actual attainment requires attainment of the finite program.
\end{itemize}

\begin{thebibliography}{99}
\bibitem{ChenGeorgiouPavon2021} Yongxin Chen, Tryphon T. Georgiou, and Michele Pavon. \emph{Stochastic Control Liaisons: Richard Sinkhorn Meets Gaspard Monge on a Schr\"odinger Bridge}. SIAM Review, 63(2):249--313, 2021. doi:10.1137/20M1339982. arXiv:2005.10963.
\bibitem{Leonard2014Survey} Christian L\'eonard. \emph{A Survey of the Schr\"odinger Problem and Some of Its Connections with Optimal Transport}. Discrete and Continuous Dynamical Systems--A, 34(4):1533--1574, 2014. doi:10.3934/dcds.2014.34.1533. arXiv:1308.0215.
\bibitem{Fortet1940} Robert Fortet. \emph{Resolution d'un systeme d'equations de M. Schrodinger}. Journal de mathematiques pures et appliquees, 9e serie, 19(1--4):83--105, 1940.
\bibitem{SinkhornKnopp1967} Richard Sinkhorn and Paul Knopp. \emph{Concerning Nonnegative Matrices and Doubly Stochastic Matrices}. Pacific Journal of Mathematics, 21(2):343--348, 1967.
\bibitem{FranklinLorenz1989} Joel Franklin and Jens Lorenz. \emph{On the Scaling of Multidimensional Matrices}. Linear Algebra and its Applications, 114/115:717--735, 1989. doi:10.1016/0024-3795(89)90490-4.
\bibitem{Carroll2004} Joseph E. Carroll. \emph{Birkhoff's Contraction Coefficient}. Linear Algebra and its Applications, 389:227--234, 2004. doi:10.1016/j.laa.2004.02.039.
\bibitem{EvesonNussbaum1995} Simon P. Eveson and Roger D. Nussbaum. \emph{An Elementary Proof of the Birkhoff--Hopf Theorem}. Mathematical Proceedings of the Cambridge Philosophical Society, 117(1):31--55, 1995.
\bibitem{BlanchetMurthy2019} Jose Blanchet and Karthyek R. A. Murthy. \emph{Quantifying Distributional Model Risk via Optimal Transport}. Mathematics of Operations Research, 44(2):565--600, 2019. doi:10.1287/moor.2018.0936. arXiv:1604.01446.
\bibitem{GaoKleywegt2023} Rui Gao and Anton J. Kleywegt. \emph{Distributionally Robust Stochastic Optimization with Wasserstein Distance}. Mathematics of Operations Research, 48(2):603--655, 2023. doi:10.1287/moor.2022.1275. arXiv:1604.02199.
\bibitem{MohajerinEsfahaniKuhn2018} Peyman Mohajerin Esfahani and Daniel Kuhn. \emph{Data-Driven Distributionally Robust Optimization Using the Wasserstein Metric: Performance Guarantees and Tractable Reformulations}. Mathematical Programming, 171(1--2):115--166, 2018. doi:10.1007/s10107-017-1172-1. arXiv:1505.05116.
\bibitem{WangGaoXie2025} Jie Wang, Rui Gao, and Yao Xie. \emph{Sinkhorn Distributionally Robust Optimization}. Operations Research, 2025. doi:10.1287/opre.2023.0294. arXiv:2109.11926.
\bibitem{Girsanov1960} I. V. Girsanov. \emph{On Transforming a Certain Class of Stochastic Processes by Absolutely Continuous Substitution of Measures}. Theory of Probability and Its Applications, 5(3):285--301, 1960. doi:10.1137/1105027.
\bibitem{CameronMartin1945} R. H. Cameron and W. T. Martin. \emph{Transformations of Wiener Integrals under a General Class of Linear Transformations}. Transactions of the American Mathematical Society, 58(2):184--219, 1945.
\bibitem{Yeh1978} J. Yeh. \emph{Transformation of Conditional Wiener Integrals under Translation and the Cameron--Martin Translation Theorem}. Tohoku Mathematical Journal, 30(4):505--515, 1978.
\bibitem{Kakutani1948} Shizuo Kakutani. \emph{On Equivalence of Infinite Product Measures}. Annals of Mathematics, 49(1):214--224, 1948.
\bibitem{ShiDeBortoliCampbellDoucet2023} Yuyang Shi, Valentin De Bortoli, Andrew Campbell, and Arnaud Doucet. \emph{Diffusion Schr\"odinger Bridge Matching}. Advances in Neural Information Processing Systems, 2023. arXiv:2303.16852.
\bibitem{LiuLipmanNickelKarrerTheodorouChen2024} Guan-Horng Liu, Yaron Lipman, Maximilian Nickel, Brian Karrer, Evangelos A. Theodorou, and Ricky T. Q. Chen. \emph{Generalized Schr\"odinger Bridge Matching}. International Conference on Learning Representations, 2024. arXiv:2310.02233.
\bibitem{DomingoEnrichDrozdzalKarrerChen2025} Carles Domingo-Enrich, Michal Drozdzal, Brian Karrer, and Ricky T. Q. Chen. \emph{Adjoint Matching: Fine-Tuning Flow and Diffusion Generative Models with Memoryless Stochastic Optimal Control}. arXiv preprint arXiv:2409.08861, 2025.

\bibitem{Birkhoff1957} Garrett Birkhoff. \emph{Extensions of Jentzsch's Theorem}. Transactions of the American Mathematical Society, 85(1):219--227, 1957.
\bibitem{Sinkhorn1964} Richard Sinkhorn. \emph{A Relationship Between Arbitrary Positive Matrices and Doubly Stochastic Matrices}. Annals of Mathematical Statistics, 35(2):876--879, 1964.
\bibitem{Sinkhorn1967Prescribed} Richard Sinkhorn. \emph{Diagonal Equivalence to Matrices with Prescribed Row and Column Sums}. American Mathematical Monthly, 74(4):402--405, 1967.
\bibitem{Bauer1965} Friedrich L. Bauer. \emph{An Elementary Proof of the Hopf Inequality for Positive Operators}. Numerische Mathematik, 7:331--337, 1965.
\bibitem{Bushell1973} P. J. Bushell. \emph{Hilbert's Metric and Positive Contraction Mappings in a Banach Space}. Archive for Rational Mechanics and Analysis, 52:330--338, 1973.

\end{thebibliography}

\end{document}
